% !TEX root = ../00_Main/Main.tex
\documentclass[../00_Main/Main.tex]{subfiles}
\begin{document}

\label{sec:methods}

Este cap\'itulo describe el entorno, la m\'etrica de desempe\~no, la
implementaci\'on de los modelos individuales y de los ensambles, la
optimizaci\'on de hiperpar\'ametros y el protocolo de evaluaci\'on.

\subsection{Implementaci\'on}

Los experimentos se implementaron en el repositorio
\texttt{mPES}~\citep{mPES2026}, que extiende el repositorio \texttt{PES}
del BCI-NE Lab~\citep{BCINE2022}, del que hereda el entorno y la
implementaci\'on tabular de Q-Learning. El c\'odigo est\'a organizado en
tres grupos de paquetes: \texttt{tabular/} (m\'etodos tabulares),
\texttt{ml/} (redes neuronales) y \texttt{ens/} (ensambles). Cada
paquete incluye el entorno, implementado con
Gymnasium~\citep{Towers2024}, su rutina de entrenamiento y, cuando
corresponde, su estudio de optimizaci\'on bayesiana. Se us\'o Python 3.12
con TensorFlow 2.21, Keras 3.13, NumPy 2.4, Optuna 4.7, Gymnasium 1.2 y
SciPy 1.17. Los comandos para reproducir cada resultado est\'an en el
Ap\'endice~\ref{ap:repro}.

\subsection{Entorno PES}

\subsubsection{Estados y acciones}

El entorno se implementa como un MDP finito. En cada paso (\emph{trial}) llega una
ciudad y el agente decide cu\'antas unidades de un recurso escaso le
asigna. El estado que observa el agente es:

\begin{equation}
\mathcal{S} \;=\;
\bigl\{(R, t, S) : R \in \{0, \ldots, 30\},\;
t \in \{0, \ldots, T_{\max}\},\;
S \in \{0, \ldots, S_{\max}\}\bigr\}
\label{eq:state-space}
\end{equation}

Donde $R$ son los recursos que quedan, $t$ el n\'umero de paso dentro
de la secuencia ($T_{\max} = 10$) y $S$ la severidad inicial de la
ciudad actual ($S_{\max} = 9$); hay $31 \times 11 \times 10 = 3\,410$
estados. Cada secuencia dispone de $39$ recursos, de los cuales el
entorno reserva $9$ como preasignados, por lo que el agente administra
$R \le 30$. Las acciones son $\mathcal{A} = \{0, 1, \ldots, 10\}$; una
asignaci\'on mayor que $R$ se recorta a los recursos disponibles.

\subsubsection{Din\'amica y recompensa}\label{sec:env-dyn}

En cada paso, la severidad de toda ciudad $i$ ya incorporada se
actualiza con:

\begin{equation}
S_{i,t+1} \;=\; \max\!\bigl(0,\ \beta\, S_{i,t} - \alpha\, a_{i}\bigr)
\qquad \alpha = 0{,}4 \quad \beta = 1 + \alpha = 1{,}4
\label{eq:transicion}
\end{equation}

Donde $a_i$ es la asignaci\'on que recibi\'o la ciudad $i$ al
incorporarse, $\alpha$ la eficacia de cada unidad de recurso y $\beta$
el factor de crecimiento por paso. La asignaci\'on de una ciudad se fija
una sola vez, pero su severidad se actualiza en todos los pasos
siguientes; por eso una asignaci\'on insuficiente se amplifica
geom\'etricamente. La recompensa de cada paso es la suma, con signo
negativo, de las severidades de las ciudades activas:

\begin{equation}
r_t \;=\; -\sum_{i \le t} S_{i,t+1}
\label{eq:reward}
\end{equation}

Maximizar el retorno equivale, por lo tanto, a minimizar la severidad
acumulada de la secuencia.

El estado que observa el agente no incluye la severidad de las ciudades
ya incorporadas ni la longitud de la secuencia. Esa severidad forma parte
de la recompensa, pero no depende de las decisiones posteriores, porque
la asignaci\'on de cada ciudad se fija al llegar. Por lo tanto no cambia
qu\'e acci\'on conviene en cada paso, aunque s\'i agrega al retorno una
parte que el estado actual no explica. En sentido estricto, $(R, t, S)$
es una observaci\'on parcial del entorno.

\subsubsection{Secuencias de evaluaci\'on}

Cada experimento se organiza en $8$ bloques de $8$ secuencias, de
$3$ a $10$ pasos cada una, lo que suma $360$ pasos. La evaluaci\'on de
referencia usa las $n = 64$ secuencias fijas de los archivos
\texttt{initial\_severity.csv} y \texttt{sequence\_lengths.csv}, id\'enticos
en todos los paquetes, de modo que todos los modelos enfrentan
exactamente las mismas secuencias. En esos archivos las severidades
iniciales van de $2$ a $8$.

El entrenamiento no recorre esas $64$ secuencias en orden: en cada
episodio sortea una longitud y unas severidades iniciales a partir de
sus distribuciones emp\'iricas en los mismos archivos. Por eso, durante
el entrenamiento las severidades s\'olo toman valores de $2$ a $8$,
aunque el espacio de estados admite de $0$ a $9$.

\subsection{M\'etrica de desempe\~no}\label{sec:metric}

La severidad final acumulada de una secuencia de $T$ pasos es:

\begin{equation}
S_{\mathrm{cruda}} = \sum_{i=1}^{T} S_i^{\mathrm{final}}
\label{eq:raw-severity-materials}
\end{equation}

Donde $S_i^{\mathrm{final}}$ es la severidad de la ciudad $i$ al
terminar la secuencia; cuanto menor, mejor. Como esta cantidad depende
de las severidades iniciales y de la longitud de cada secuencia, se
normaliza con dos referencias: $S_{\mathrm{peor}}$, la severidad
acumulada si no se asigna ning\'un recurso, y $S_{\mathrm{mejor}}$, la
m\'inima severidad acumulada que se puede lograr con el presupuesto de
la secuencia conociendo de antemano todas sus ciudades. Como la
asignaci\'on de cada ciudad se fija al llegar y su evoluci\'on no depende
de las dem\'as, $S_{\mathrm{mejor}}$ se calcula de forma exacta por
programaci\'on din\'amica, como un problema de mochila acotada sobre
asignaciones enteras de $0$ a $10$ por ciudad que suman a lo sumo $30$.
Ambas referencias se calculan al terminar la secuencia y no intervienen
en las decisiones del agente. El \emph{desempe\~no normalizado} de una
secuencia es:

\begin{equation}
\bar{r} =
\frac{S_{\mathrm{peor}} - S_{\mathrm{agente}}}
{S_{\mathrm{peor}} - S_{\mathrm{mejor}}}
\label{eq:normalised-severity-materials}
\end{equation}

Donde $S_{\mathrm{agente}}$ es la severidad cruda que obtuvo el agente.
As\'i, $\bar{r}=0$ equivale a no asignar recursos y $\bar{r}=1$ a la
asignaci\'on \'optima; como ning\'un agente puede gastar m\'as que el
presupuesto, $\bar r \le 1$. Como referencia se evalu\'o un agente que
elige acciones al azar sobre las mismas $64$ secuencias
(Figura~\ref{fig:baseline-random}):
obtiene $\bar{r} = 0{,}670$ en promedio (desviaci\'on est\'andar
$0{,}186$, rango $[0{,}164;\,0{,}950]$). Gastar el presupuesto sin
criterio reduce parte de la severidad, pero queda lejos del \'optimo en
la mayor\'ia de las secuencias.

\begin{figure}[!htb]
  \centering
  \begin{subfigure}[t]{0.49\textwidth}
    \centering
    \includegraphics[width=\linewidth]{random_player_sequence_performance.png}
    \caption[Severidad cruda]{Severidad final acumulada $S_{\mathrm{cruda}}$.}
    \label{fig:baseline-random-raw}
  \end{subfigure}
  \hfill
  \begin{subfigure}[t]{0.49\textwidth}
    \centering
    \includegraphics[width=\linewidth]{random_player_normalised_performance.png}
    \caption[Desempe\~no normalizado]{Desempe\~no normalizado $\bar{r}$;
    las l\'ineas horizontales marcan $\bar{r}=0$ y $\bar{r}=1$.}
    \label{fig:baseline-random-normalised}
  \end{subfigure}
  \caption[Agente aleatorio]{Agente aleatorio (\emph{decisor aleatorio}
  en la figura) en las $64$ secuencias de evaluaci\'on. La l\'inea
  punteada indica la media de cada serie.}
  \label{fig:baseline-random}
\end{figure}

\subsection{Modelos individuales}

La Tabla~\ref{tab:packages} lista los modelos individuales, el paquete
que los implementa y el nombre con que se los menciona en el resto del
texto; los ensambles se listan en la Tabla~\ref{tab:ens-summary}. Las
figuras generadas por el c\'odigo usan el nombre del paquete.

\begin{table}[!htb]
\centering
\caption[Modelos individuales]{Modelos individuales, paquete que los
implementa y nombre usado en el texto. Q-Learning base se incluye s\'olo
como referencia sin optimizaci\'on.}
\label{tab:packages}
\small
\begin{tabular}{@{}lll@{}}
\toprule
\textbf{Paquete} & \textbf{Nombre en el texto} & \textbf{Descripci\'on} \\
\midrule
\texttt{pes\_base} & Q-Learning base & Q-Learning tabular sin optimizar.\\
\texttt{pes\_ql}   & Q-Learning & Q-Learning tabular optimizado.\\
\texttt{pes\_dql}  & Double Q-Learning & Double Q-Learning tabular con PBRS.\\
\texttt{pes\_dqn}  & DQN & Double DQN con red densa.\\
\texttt{pes\_rdqn} & DQN recurrente & Double DQN con codificador LSTM.\\
\texttt{pes\_a2c}  & A2C & Actor--cr\'itico.\\
\texttt{pes\_trf}  & Transformer & Double DQN con codificador Transformer causal.\\
\bottomrule
\end{tabular}
\end{table}

\subsubsection{Modelos tabulares}\label{sec:tabular}

Los tres modelos tabulares guardan $Q(s,a)$ en una tabla de
$31\times11\times10\times11=37\,510$ valores, indexada por los enteros
$(R, t, S)$ y la acci\'on. Las tablas se inicializan con valores de una
distribuci\'on uniforme $\mathcal{U}(-1,1)$, las acciones que exceden los
recursos se enmascaran y, durante el entrenamiento, se elige la acci\'on
con una pol\'itica $\varepsilon$-\emph{greedy}. La actualizaci\'on de
Q-Learning es:

\begin{equation}
Q(s,a) \leftarrow Q(s,a) + \alpha
\left[r + \gamma\max_{a'}Q(s',a') - Q(s,a)\right]
\label{eq:tabular-update-materials}
\end{equation}

Salvo en el estado terminal, donde $Q(s,a)\leftarrow r$. En la
evaluaci\'on se elige la acci\'on factible de mayor valor. La tabla s\'olo
tiene filas para severidades de $0$ a $9$: si la severidad es mayor, el
\'indice se recorta a $9$. Como el entrenamiento s\'olo presenta
severidades de $2$ a $8$, las filas de $S = 0$, $1$ y $9$ nunca se
actualizan y conservan sus valores iniciales aleatorios en las tres
tablas entrenadas.

\paragraph{Q-Learning base y Q-Learning} Ambos aplican la
Ecuaci\'on~\eqref{eq:tabular-update-materials} con una exploraci\'on
$\varepsilon$ que decrece linealmente desde $\varepsilon_0$ hasta
$\varepsilon_{\min}$. Q-Learning base usa los valores originales del
repositorio PES; Q-Learning usa los hiperpar\'ametros elegidos por
optimizaci\'on bayesiana (Tabla~\ref{tab:tabular-hparams}).

\paragraph{Double Q-Learning} Mantiene dos tablas $Q_A$ y $Q_B$
(Ecuaci\'on~\eqref{eq:double-q}). En cada paso se actualiza una sola,
elegida con probabilidad $1/2$, y la pol\'itica final es
$\arg\max_a[Q_A(s,a)+Q_B(s,a)]$. La exploraci\'on se mantiene en
$\varepsilon_0$ durante un calentamiento y luego decae
exponencialmente:

\begin{equation}
\lambda = \left(\frac{\varepsilon_{\min}}{\varepsilon_0}\right)^{1/((q-w)N)}
\qquad
\varepsilon_t = \max\!\left(\varepsilon_{\min},\;
\varepsilon_0\,\lambda^{\,t-wN}\right) \quad \text{para } t \ge wN
\label{eq:dql-epsilon-materials}
\end{equation}

Donde $t$ es el n\'umero de episodio, $N$ el total de episodios, $w$ la
fracci\'on de episodios de calentamiento, $q$ la fracci\'on en la que se
alcanza $\varepsilon_{\min}$ y $\lambda$ el factor de decaimiento por
episodio. Adem\'as usa PBRS con $\Phi(s)=-\sum_i S_i$, la suma de las
severidades actuales de las ciudades con signo negativo, de modo que la
recompensa con que se actualiza la tabla es:

\begin{equation}
\tilde r = r + \kappa\,[\gamma\Phi(s')-\Phi(s)]
\label{eq:dql-pbrs-materials}
\end{equation}

Donde $\kappa$ es el coeficiente de PBRS y $s'$ el estado siguiente.

\begin{table}[!htb]
\centering
\caption[Hiperpar\'ametros de los modelos tabulares]{Hiperpar\'ametros
de los modelos tabulares. Los de Q-Learning y Double Q-Learning son los
del mejor ensayo de su optimizaci\'on bayesiana.}
\label{tab:tabular-hparams}
\small
\begin{tabular}{@{}lccc@{}}
\toprule
 & \textbf{Q-Learning base} & \textbf{Q-Learning} & \textbf{Double Q-Learning} \\
\midrule
Tasa de aprendizaje $\alpha$ & $0{,}2$ & $0{,}2862$ & $0{,}1132$ \\
Descuento $\gamma$ & $0{,}9$ & $0{,}8587$ & $0{,}9777$ \\
$\varepsilon_0$ / $\varepsilon_{\min}$ & $0{,}8$ / $0$ & $0{,}6813$ / $0{,}0437$ & $0{,}5998$ / $0{,}0327$ \\
Decaimiento de $\varepsilon$ & lineal & lineal & exponencial \\
Calentamiento $w$ & --- & --- & $0{,}0260$ \\
Fracci\'on $q$ & --- & --- & $0{,}6430$ \\
PBRS $\kappa$ & --- & --- & $0{,}000392$ \\
Episodios & $1\,000\,000$ & $550\,000$ & $360\,000$ \\
Ensayos de optimizaci\'on & --- & $100$ & $100$ \\
\bottomrule
\end{tabular}
\end{table}

\subsubsection{Modelos profundos}\label{sec:deep}

\paragraph{Entrada y configuraci\'on com\'un} Los cuatro modelos
profundos reciben el estado escalado por las cotas del entorno,
$s = [R/30,\; t/10,\; S/9] \in [0,1]^{3}$, para que las tres componentes
tengan rangos comparables en la entrada de la red. Es una funci\'on fija
de las cotas del MDP, que se aplica igual en entrenamiento,
evaluaci\'on y ensambles. Con severidades mayores que $9$, la entrada
$S/9$ de DQN, DQN recurrente y Transformer simplemente supera $1$; la
evaluaci\'on de A2C, en cambio, recorta la severidad a $9$ antes de
escalarla. Todos usan el optimizador
Adam~\citep{Kingma2015}, inicializaci\'on Glorot
uniforme~\citep{Glorot2010} y recorte de la norma del gradiente. La
semilla de entrenamiento de cada modelo es $42 + i + 1$, con $i$ el
\'indice del mejor ensayo de su optimizaci\'on.

DQN, DQN recurrente y Transformer comparten adem\'as la regla de aprendizaje. Usan el blanco Double DQN (Ecuaci\'on~\eqref{eq:ddqn-target}) con enmascarado de acciones no factibles y p\'erdida de Huber, y exploran con una pol\'itica $\varepsilon$-\emph{greedy} seg\'un el esquema de la Ecuaci\'on~\eqref{eq:dql-epsilon-materials}. Empiezan a aprender cuando el \emph{buffer} alcanza una fracci\'on $f$ de su capacidad y usan PBRS cuando $\kappa > 0$. Se diferencian en c\'omo procesan la entrada y en
sus hiperpar\'ametros (Tabla~\ref{tab:deep-hparams}).

\paragraph{DQN} Red densa
\texttt{Input}$(3) \!\to\! \texttt{Dense}(64,\!\mathrm{ReLU}) \!\to\!
\texttt{Dense}(64,\!\mathrm{ReLU}) \!\to\! \texttt{Dense}(11)$, que
produce un valor $Q$ por acci\'on a partir del estado actual,
con repetici\'on de experiencia~\citep{Lin1992} y red
objetivo~\citep{Mnih2015}.

\paragraph{DQN recurrente} Sigue el esquema DRQN~\citep{Hausknecht2015}.
La red recibe los \'ultimos $W = 6$ estados escalados
$(s_{t-5}, \ldots, s_t)$, completados con ceros al principio de cada
secuencia. Una LSTM de $64$ unidades~\citep{Hochreiter1997} recorre la
ventana y su \'ultimo estado oculto $h_t$ (Ecuaci\'on~\eqref{eq:lstm})
pasa por dos capas densas de $96$ unidades con ReLU y una salida de
$11$ valores $Q$. Cada ventana se procesa de forma independiente: el
estado oculto no se conserva entre decisiones y el \emph{buffer}
guarda la ventana completa de cada transici\'on. La b\'usqueda bayesiana
incluy\'o $W$, el ancho de la LSTM y el de las capas densas, pero el
entrenamiento final s\'olo toma del mejor ensayo las capas densas y los
hiperpar\'ametros de entrenamiento: $W = 6$ y las $64$ unidades de la
LSTM son los valores fijos de la configuraci\'on del paquete (el mejor
ensayo usaba $32$ unidades).

\paragraph{Transformer} Recibe la misma ventana de $W = 6$ estados. Cada
estado se proyecta a dimensi\'on $d_{\mathrm{model}} = 32$ y se le suma
un vector de posici\'on fijo, inicializado con Glorot uniforme a partir
de la semilla y no entrenado. Cada uno de los $L = 2$ bloques tiene dos
subcapas: atenci\'on causal (Ecuaci\'on~\eqref{eq:causal-attention}) con
$h = 4$ cabezas de dimensi\'on $16$ y una capa \emph{feed-forward} de
$64$ unidades. Cada subcapa tiene conexi\'on residual y va precedida de
normalizaci\'on de capa (disposici\'on \emph{Pre-LN}~\citep{Parisotto2020}).
No se usa \emph{dropout}. La representaci\'on de la \'ultima
posici\'on pasa por una capa densa de $32$ unidades con ReLU y una
salida de $11$ valores $Q$. La b\'usqueda bayesiana incluy\'o la ventana
$W$ y la arquitectura del codificador ($d_{\mathrm{model}}$, n\'umero de
cabezas y de bloques, dimensi\'on de las claves, ancho de la capa
\emph{feed-forward} y \emph{dropout}), pero el entrenamiento final s\'olo
toma del mejor ensayo la capa densa de salida y los hiperpar\'ametros de
entrenamiento. La ventana y el codificador descritos son los valores
fijos de la configuraci\'on del paquete; el mejor ensayo usaba $W = 3$,
$d_{\mathrm{model}} = 16$, $8$ cabezas, $4$ bloques, una capa
\emph{feed-forward} de $128$ unidades y \emph{dropout} de $0{,}15$.

\begin{table}[!htb]
\centering
\caption[Hiperpar\'ametros de los modelos de tipo DQN]{Arquitectura e
hiperpar\'ametros de DQN, DQN recurrente y Transformer, tal como se
entrenaron los modelos desplegados. Entre par\'entesis, el n\'umero de
transiciones que corresponde a la fracci\'on $f$. En DQN, todas las filas
vienen del mejor ensayo de la optimizaci\'on bayesiana. En DQN recurrente
y Transformer, la entrada, los par\'ametros entrenables y los ensayos de
optimizaci\'on corresponden a la arquitectura fija de \texttt{config/CONFIG.py}
del paquete; el resto de las filas (tasa de aprendizaje, $\gamma$,
$\varepsilon$, fracciones $w$/$q$, lote, \emph{buffer}, sincronizaci\'on,
recorte del gradiente, PBRS, inicio del aprendizaje y episodios) s\'i
vienen del mejor ensayo de cada paquete (v\'ease \S\ref{sec:deep} para la
discrepancia entre la arquitectura del mejor ensayo y la desplegada).}
\label{tab:deep-hparams}
\small
\begin{tabular}{@{}lccc@{}}
\toprule
 & \textbf{DQN} & \textbf{DQN recurrente} & \textbf{Transformer} \\
\midrule
Entrada & estado actual & ventana de $6$ estados & ventana de $6$ estados \\
Par\'ametros entrenables & $5\,131$ & $34\,027$ & $27\,019$ \\
Tasa de aprendizaje & $0{,}001508$ & $0{,}002415$ & $0{,}000215$ \\
Descuento $\gamma$ & $0{,}9634$ & $0{,}9719$ & $0{,}9234$ \\
$\varepsilon_0$ / $\varepsilon_{\min}$ & $0{,}9627$ / $0{,}0691$ & $0{,}8883$ / $0{,}0887$ & $0{,}8651$ / $0{,}0838$ \\
Fracciones $w$ / $q$ & $0{,}2779$ / $0{,}6290$ & $0{,}1237$ / $0{,}5990$ & $0{,}2428$ / $0{,}5333$ \\
Tama\~no de lote & $128$ & $64$ & $128$ \\
Capacidad del \emph{buffer} & $20\,000$ & $60\,000$ & $30\,000$ \\
Sincronizaci\'on $C$ (pasos) & $1\,000$ & $2\,000$ & $500$ \\
Recorte del gradiente & $3{,}953$ & $3{,}475$ & $4{,}170$ \\
PBRS $\kappa$ & $0{,}0226$ & $0{,}00263$ & $0$ \\
Inicio del aprendizaje $f$ & $0{,}1615$ ($3\,230$) & $0{,}2021$ ($12\,128$) & $0{,}1217$ ($3\,650$) \\
Episodios & $40\,000$ & $30\,000$ & $30\,000$ \\
Ensayos de optimizaci\'on & $47$ & $20$ & $10$ \\
\bottomrule
\end{tabular}
\end{table}

\paragraph{A2C} Actor y cr\'itico son redes separadas:
\texttt{Input}$(3) \!\to\! \texttt{Dense}(128,\!\mathrm{ReLU})
\!\to\! \texttt{Dense}(11,\!\mathrm{softmax})$ para el actor, que produce
directamente una distribuci\'on sobre las $11$ acciones, e
\texttt{Input}$(3) \!\to\! \texttt{Dense}(128,\!\mathrm{ReLU}) \!\to\!
\texttt{Dense}(1)$ para el cr\'itico. La ventaja se estima con GAE
(Ecuaci\'on~\eqref{eq:gae}) y se estandariza por lote; el gradiente de
pol\'itica (Ecuaci\'on~\eqref{eq:pg}) y la entrop\'ia se calculan sobre la
distribuci\'on restringida a las acciones factibles~\citep{Mnih2016}.
La Tabla~\ref{tab:a2c-hparams} resume sus hiperpar\'ametros.

\begin{table}[!htb]
\centering
\caption[Hiperpar\'ametros de A2C]{Hiperpar\'ametros de A2C (mejor ensayo
de $100$ de la optimizaci\'on bayesiana).}
\label{tab:a2c-hparams}
\small
\begin{tabular}{@{}ll@{}}
\toprule
Tasa de aprendizaje inicial (actor / cr\'itico) & $0{,}000648$ / $0{,}004299$ \\
Decaimiento de la tasa & coseno, hasta el $23{,}75\,\%$ del valor inicial \\
Descuento $\gamma$ & $0{,}8545$ \\
Coeficiente de entrop\'ia $\beta_H$ & $0{,}00528$ \\
$\lambda$ de GAE & $0{,}9135$ \\
Recorte del gradiente (ambas redes) & $1{,}200$ \\
PBRS $\kappa$ & $0{,}1527$ \\
Penalizaci\'on de gasto & $r \leftarrow r - 0{,}01039\,a$ \\
Sesgo inicial del logit de $a = 10$ & $-1{,}389$ \\
Episodios & $125\,000$ \\
\bottomrule
\end{tabular}
\end{table}

\subsection{Ensambles}\label{sec:ens-methods}

Los ensambles no se entrenan: cargan los modelos ya entrenados de DQN,
DQN recurrente y Transformer y combinan sus salidas en cada decisi\'on.
Cinco de las seis variantes suman el actor de A2C como cuarto miembro.
La Tabla~\ref{tab:ens-summary} resume qu\'e hace cada regla; el resto de
la secci\'on las define con precisi\'on, con una notaci\'on com\'un.

\begin{table}[!htb]
\centering
\caption[Ensambles evaluados]{Ensambles evaluados, paquete que los
implementa y qu\'e hace cada regla de combinaci\'on.}
\label{tab:ens-summary}
\small
\begin{tabular}{@{}>{\raggedright\arraybackslash}p{0.19\linewidth}l>{\raggedright\arraybackslash}p{0.44\linewidth}@{}}
\toprule
\textbf{Ensamble} & \textbf{Paquete} & \textbf{C\'omo decide} \\
\midrule
Ensamble ponderado & \texttt{pes\_ens} & Promedia las distribuciones de DQN,
DQN recurrente y Transformer con pesos que crecen con la confianza; el
Transformer tiene el mayor peso. Agrega un \emph{prior} de severidad y
una cota de seguridad.\\
Voto suave & \texttt{pes\_ens\_sprb} & Promedia las distribuciones de los
cuatro miembros, ponderadas por peso y confianza.\\
Voto por acci\'on & \texttt{pes\_ens\_accq} & Cada miembro vota por su
acci\'on preferida; los empates se resuelven con los valores $Q$.\\
Consenso & \texttt{pes\_ens\_consensus} & Voto ponderado m\'as un premio
por acuerdo y una penalizaci\'on por desacuerdo entre miembros
confiados.\\
Consenso con prior & \texttt{pes\_ens\_consensus\_prior} & Como el
consenso, pero sumando distribuciones y mezclando el resultado con un
\emph{prior} de severidad.\\
Compuerta del Transformer & \texttt{pes\_ens\_trf\_guard} & Sigue al
Transformer salvo cuando su confianza es baja; entonces usa un voto suave
de respaldo.\\
\bottomrule
\end{tabular}
\end{table}

\subsubsection{Notaci\'on com\'un}

En el estado $s$, con $R$ recursos disponibles, para cada miembro $k$ se
calculan:

\begin{itemize}
    \item $p_k(a \mid s)$: su distribuci\'on sobre las acciones factibles
    $a \le R$. Para DQN, DQN recurrente y Transformer se obtiene
    aplicando a sus valores $Q_k(s,a)$ la \emph{softmax} de
    temperatura $\tau$ (Ecuaci\'on~\eqref{eq:ens-softmax}) y
    renormalizando sobre las acciones factibles; para A2C es la salida
    del actor, renormalizada del mismo modo.
    \item $a_k = \arg\max_a p_k(a \mid s)$: la acci\'on que prefiere.
    \item $c_k = 1 - H_{\mathrm{norm}}(p_k)$: su confianza en la decisi\'on
    (Ecuaci\'on~\eqref{eq:shannon-norm}), entre $0$ y $1$. En esta
    secci\'on la entrop\'ia se normaliza por el logaritmo del n\'umero de
    acciones factibles y no por $\log|\mathcal{A}|$; la excepci\'on es el
    ensamble ponderado, que normaliza por $\log 11$.
    \item $w_k$: su peso base, un par\'ametro fijo de cada ensamble.
    \item $\tilde w_k = w_k\, c_k^{\rho}$: su peso efectivo, que combina el
    peso base con la confianza; el exponente $\rho \ge 0$ regula cu\'anto
    influye la confianza ($\rho = 0$ la ignora).
    \item $\hat Q_k(a) = \bigl(Q_k(a)-\overline{Q}_k\bigr)/\mathrm{sd}(Q_k)$:
    sus valores $Q$ estandarizados, con $\overline{Q}_k$ y
    $\mathrm{sd}(Q_k)$ la media y la desviaci\'on est\'andar de sus valores
    sobre las acciones factibles en el estado $s$ (para A2C se
    estandarizan sus probabilidades). Permiten comparar modelos cuyos
    valores tienen escalas distintas.
\end{itemize}

Los par\'ametros de cada ensamble se listan en la
Tabla~\ref{tab:ens-params}. Los del ensamble ponderado son valores fijos
de su configuraci\'on. Los de los otros cinco provienen de estudios de
Optuna (TPE) sobre las $64$ secuencias de evaluaci\'on. S\'olo en el
consenso con \emph{prior} el estudio guardado corresponde a la
configuraci\'on evaluada: su puntaje registrado ($0{,}918$) coincide con
la media que obtiene el ensamble en la referencia. En los otros cuatro,
el archivo de par\'ametros no incluye el peso de A2C, que toma el valor
por defecto $0{,}10$, y el puntaje registrado por el estudio (entre
$0{,}888$ y $0{,}911$) no coincide con la media evaluada
(Tabla~\ref{tab:ensemble-stress}). Esos estudios se hicieron, por lo
tanto, con una configuraci\'on de miembros distinta de la evaluada.

\subsubsection{Ensamble ponderado}

Combina DQN, DQN recurrente y Transformer ($K = 3$) con un voto suave
cuyos pesos crecen con la confianza de cada miembro:

\begin{equation}
  p_{\mathrm{ens}}(a \mid s) \;=\;
  \frac{\eta(a)\,\sum_{k=1}^{K} w_k\,(0{,}1 + c_k)\, p_k(a \mid s)}
       {\sum_{a' \le R} \eta(a')\,\sum_{k=1}^{K} w_k\,(0{,}1 + c_k)\, p_k(a' \mid s)}
  \qquad
  \eta(a) = \begin{cases} 0{,}3 & a = 0,\ R > 0 \\ 1 & \text{en otro caso} \end{cases}
  \label{eq:ens-soft}
\end{equation}

Donde $w_k$ son los pesos base ($0{,}18$, $0{,}90$ y $5{,}0$ para DQN,
DQN recurrente y Transformer, normalizados para que sumen $1$). El t\'ermino $0{,}1 + c_k$ hace que el peso crezca con la confianza sin
anularse cuando $c_k = 0$, y $\eta(a)$ reduce la probabilidad de no
asignar nada mientras queden recursos. El denominador normaliza la
distribuci\'on sobre las acciones factibles. Aqu\'i la \emph{softmax} usa
$\tau = 15$. Luego esa distribuci\'on se mezcla con una distribuci\'on
previa (\emph{prior}) centrada en la severidad actual:

\begin{equation}
  \pi_{\mathrm{prior}}(a\mid s) \propto
  \exp\!\left[-\frac{(a-S)^2}{2\sigma_S^2}\right]
  \qquad
  p'(a\mid s) \propto (1-w_\pi)\,p_{\mathrm{ens}}(a\mid s) + w_\pi\,\pi_{\mathrm{prior}}(a\mid s)
  \label{eq:ens-prior}
\end{equation}

Donde $S$ es la severidad de la ciudad actual, $\sigma_S = 3$ el ancho
del \emph{prior} y $w_\pi = 0{,}17$ su peso en la mezcla; $p'$ se
renormaliza sobre las acciones factibles. La acci\'on ejecutada es la
de mayor probabilidad:

\begin{equation}
  a_{\mathrm{ens}} \;=\; \arg\max_{a \le R}\; p'(a \mid s)
  \label{eq:ens-action}
\end{equation}

Por \'ultimo, una cota de seguridad reemplaza esa acci\'on por
$\lfloor S/2 \rfloor$ cuando $S \geq 6$, la acci\'on elegida es menor que
esa cota y la cota es factible. Como en A2C, la severidad de entrada se
recorta a $9$ antes de escalarla, tanto para los miembros como para el
\emph{prior}.

\subsubsection{Voto suave}

Promedia las distribuciones de los miembros con sus pesos efectivos y
elige la acci\'on m\'as probable:

\begin{equation}
  p_{\mathrm{vs}}(a\mid s) =
  \frac{\sum_{k} \tilde w_k\, p_k(a\mid s)}{\sum_{k} \tilde w_k}
  \qquad
  a_{\mathrm{vs}} = \arg\max_a p_{\mathrm{vs}}(a\mid s)
  \label{eq:ens-sprb}
\end{equation}

Donde la suma recorre los cuatro miembros. Es la \'unica variante que
ajusta la temperatura de la \emph{softmax} ($\tau = 0{,}1044$); en las
dem\'as $\tau = 1$.

\subsubsection{Voto por acci\'on}

Cada miembro vota por su acci\'on preferida con su peso efectivo; si hay
empate, gana la acci\'on con mayor suma de valores estandarizados:

\begin{equation}
  V(a) = \sum_{k} \tilde w_k\, \mathbf{1}[a_k=a]
  \qquad
  T(a) = \sum_{k} \hat Q_k(a)
  \label{eq:ens-accq-vote}
\end{equation}

\begin{equation}
  a_{\mathrm{va}} =
  \arg\max_{a \in \mathcal{V}} T(a)
  \qquad
  \mathcal{V} = \{a : V(a)=\textstyle\max_{a'} V(a')\}
  \label{eq:ens-accq-action}
\end{equation}

Donde $V(a)$ es el total de votos ponderados que recibe la acci\'on $a$,
$T(a)$ el criterio de desempate y $\mathcal{V}$ el conjunto de acciones
m\'as votadas.

\subsubsection{Consenso}

Parte del mismo voto ponderado y le suma un premio por acuerdo, una
penalizaci\'on por desacuerdo y los valores estandarizados de los
miembros ponderados por su confianza:

\begin{equation}
  \begin{gathered}
  G(a) ={} \sum_{k} \tilde w_k\, \mathbf{1}[a_k=a]
  \;+\; \beta_a \!\!\sum_{k:\,a_k=a} c_k
  \;-\; \beta_d \!\!\sum_{k:\,a_k\neq a} c_k
  \;+\; \sum_{k} \hat Q_k(a)\, c_k \\
  a_{\mathrm{cons}} ={} \arg\max_a G(a)
  \end{gathered}
  \label{eq:ens-consensus}
\end{equation}

Donde $G(a)$ es el puntaje de la acci\'on $a$. El segundo t\'ermino suma
la confianza de los miembros que prefieren $a$, multiplicada por el premio $\beta_a$. El tercero resta la confianza de los que prefieren otra
acci\'on, multiplicada por la penalizaci\'on $\beta_d$. El cuarto agrega
lo que cada miembro opina de $a$ aunque no sea su preferida.

\subsubsection{Consenso con \emph{prior}}

Reemplaza el voto duro del consenso por una suma de distribuciones,
\mbox{$\tilde G(a) = \sum_k \tilde w_k\, p_k(a\mid s)$}, a la que suma los
mismos t\'erminos de acuerdo y desacuerdo de la
Ecuaci\'on~\eqref{eq:ens-consensus}. Luego multiplica $\tilde G(0)$ por
$0{,}3$ cuando quedan recursos, normaliza $\tilde G$ como una
distribuci\'on y la mezcla con el \emph{prior} de severidad de la
Ecuaci\'on~\eqref{eq:ens-prior}, con sus propios $w_\pi$ y $\sigma_S$. La
acci\'on es la de mayor valor en esa mezcla, con la misma cota de
seguridad que el ensamble ponderado. A diferencia de los dem\'as, calcula
la confianza aplicando la \emph{softmax} a $p_k$ en lugar de a los
valores $Q$, lo que acerca la distribuci\'on a la uniforme y reduce
$c_k$. Es la \'unica variante en la que se
optimiz\'o el peso del actor de A2C; en las dem\'as vale $0{,}10$.

\subsubsection{Compuerta del Transformer}

Sigue al Transformer salvo cuando su confianza es baja. Para decidirlo
calcula una compuerta log\'istica:

\begin{equation}
  g = \frac{1}{1+\mathrm{e}^{-\kappa_g\,(c_{\mathrm{trf}}-\tau_g)}}
  \qquad
  f(a) = \sum_{k} \tilde w_k\, p_k(a\mid s)
  \label{eq:ens-guard}
\end{equation}

Donde $c_{\mathrm{trf}}$ es la confianza del Transformer, $\tau_g$ el
umbral de confianza, $\kappa_g$ la pendiente de la compuerta y $f(a)$ un
voto suave de respaldo con todos los miembros. Si $g \geq 0{,}5$ se
ejecuta la acci\'on del Transformer, $a_{\mathrm{trf}}$; si no, la de
mayor $f(a)$. Como $g \geq 0{,}5$ equivale a $c_{\mathrm{trf}} \ge
\tau_g$, el ensamble sigue al Transformer siempre que su confianza
supere $\tau_g = 0{,}2248$.

\begin{table}[!htb]
\centering
\caption[Par\'ametros de los ensambles]{Par\'ametros de los seis
ensambles: temperatura $\tau$, exponente de confianza $\rho$ y pesos
base de DQN ($w_{\mathrm{DQN}}$), DQN recurrente ($w_{\mathrm{rec}}$),
Transformer ($w_{\mathrm{trf}}$) y A2C ($w_{\mathrm{A2C}}$). En el
ensamble ponderado los pesos se multiplican por $0{,}1 + c_k$ en lugar de
$c_k^{\rho}$ y el actor de A2C no participa.}
\label{tab:ens-params}
\footnotesize
\setlength{\tabcolsep}{4pt}
\begin{tabular}{@{}lcccccc>{\raggedright\arraybackslash}p{3.4cm}@{}}
\toprule
\textbf{Ensamble} & $\tau$ & $\rho$ & $w_{\mathrm{DQN}}$ & $w_{\mathrm{rec}}$ & $w_{\mathrm{trf}}$ & $w_{\mathrm{A2C}}$ & \textbf{Otros} \\
\midrule
Ponderado & $15$ & --- & $0{,}18$ & $0{,}90$ & $5{,}0$ & --- & $w_\pi = 0{,}17$, $\sigma_S = 3$ \\
Voto suave & $0{,}1044$ & $0{,}4677$ & $2{,}852$ & $1{,}628$ & $1{,}134$ & $0{,}10$ & --- \\
Voto por acci\'on & $1$ & $1{,}005$ & $2{,}374$ & $0{,}722$ & $1{,}008$ & $0{,}10$ & --- \\
Consenso & $1$ & $0{,}4026$ & $2{,}995$ & $0{,}320$ & $1{,}695$ & $0{,}10$ & $\beta_a = 0{,}1649$, $\beta_d = 0{,}3647$ \\
Consenso con prior & $1$ & $0{,}4180$ & $2{,}778$ & $2{,}465$ & $2{,}840$ & $0{,}1044$ & $\beta_a = 1{,}810$, $\beta_d = 0{,}3462$, $w_\pi = 0{,}2051$, $\sigma_S = 0{,}9177$ \\
Compuerta Transformer & $1$ & $1{,}896$ & $1{,}124$ & $2{,}852$ & $2{,}196$ & $0{,}10$ & $\tau_g = 0{,}2248$, $\kappa_g = 3{,}964$ \\
\bottomrule
\end{tabular}
\end{table}

Como las reglas no se entrenan junto con los miembros, las diferencias
entre ensambles miden el efecto de la regla de combinaci\'on y de sus
par\'ametros sobre los mismos modelos ya entrenados.

\subsection{Optimizaci\'on de hiperpar\'ametros}

Cada paquete optimizable tiene un m\'odulo \texttt{optimize\_*.py} que
ejecuta un estudio de Optuna con TPE~\citep{Bergstra2011,Akiba2019}. En
cada ensayo se entrena el agente desde cero con una configuraci\'on
$\boldsymbol{\theta}$ y se calcula el desempe\~no medio $\bar{r}$
(Ecuaci\'on~\eqref{eq:normalised-severity-materials}) sobre las $64$
secuencias; el estudio busca
$\boldsymbol{\theta}^{*}=\arg\max_{\boldsymbol{\theta}} \bar r(\boldsymbol{\theta})$.
Los ensayos poco prometedores se interrumpen con \texttt{MedianPruner}
y el estudio se guarda en una base SQLite, lo que permite reanudarlo. El
modelo final se reentrena desde cero con $\boldsymbol{\theta}^{*}$ y la
semilla de su ensayo. El n\'umero de episodios forma parte del espacio de
b\'usqueda, y el entrenamiento final usa el del mejor ensayo. Los n\'umeros
de ensayos y de episodios de cada modelo est\'an en las
Tablas~\ref{tab:tabular-hparams}, \ref{tab:deep-hparams}
y~\ref{tab:a2c-hparams}; A2C us\'o $100$ ensayos. En Double Q-Learning,
$86$ de los $100$ ensayos se interrumpieron por el \texttt{MedianPruner}.

En Q-Learning, Double Q-Learning, DQN y A2C, el modelo final reproduce
el puntaje de su mejor ensayo. En el DQN recurrente y el Transformer, en
cambio, el reentrenamiento conserva la arquitectura fija del paquete
y no la del mejor ensayo. El DQN recurrente final obtiene $0{,}899$, por
debajo de los $0{,}926$ de su mejor ensayo; el Transformer final obtiene
$0{,}927$, frente a $0{,}925$.

\subsection{Escenarios de generalizaci\'on}\label{sec:scenario-catalogue}

Adem\'as de la condici\'on de referencia (\texttt{sev\_base}), que
reproduce la distribuci\'on de entrenamiento, cada modelo se eval\'ua en
$21$ escenarios de generalizaci\'on (Tabla~\ref{tab:scenarios}),
agrupados en cuatro familias:

\begin{itemize}
  \item \textbf{Severidad ($9$).} Cambian la distribuci\'on de las
  severidades iniciales (uniforme, gaussianas de media baja, media y
  alta, Weibull, Beta sesgadas y bimodal). Uno de ellos, la
  \emph{extrapolaci\'on de severidad} (\texttt{sev\_extrapolate\_high}),
  usa severidades de $10$ a $12$, por encima de la cota del entorno
  ($S_{\max} = 9$). Los otros ocho incluyen severidades de $0$, $1$ o
  $9$, que est\'an dentro del espacio de estados pero no aparecen en el
  entrenamiento ($2$ a $8$).
  \item \textbf{Longitud ($5$).} Cambian la longitud de las secuencias
  (todas de $3$, todas de $10$, geom\'etrica, Poisson). Uno de ellos, el
  de \emph{secuencias largas} (\texttt{len\_extrapolate\_long}), usa
  secuencias de $11$ a $20$ pasos, por encima de la cota del entorno
  ($T_{\max} = 10$).
  \item \textbf{Conjunta ($4$).} Cambian severidad y longitud a la vez;
  la \emph{extrapolaci\'on conjunta} (\texttt{joint\_extrap\_both})
  lleva ambas por encima de sus cotas.
  \item \textbf{Estructural ($3$).} Cambian s\'olo la partici\'on en
  bloques o el n\'umero de secuencias ($n = 128$ en uno de ellos), sin
  cambiar la distribuci\'on. Sirven de control: deben reproducir el
  resultado de la referencia.
\end{itemize}

En el resto del texto se usan los nombres en cursiva; los
identificadores aparecen en la Tabla~\ref{tab:scenarios} y en las
figuras. Se llama \emph{fuera de rango} a los tres escenarios que superan
las cotas del entorno.

Los escenarios de generalizaci\'on se generan con la semilla $42$, de
modo que, dentro de cada escenario, todos los modelos enfrentan las
mismas secuencias. La
evaluaci\'on es s\'olo de inferencia: ning\'un modelo se reentrena. Para
cada modelo y escenario se calcula la \emph{degradaci\'on}
$\Delta\bar{r} = \bar{r}_{\text{ref}} - \bar{r}_{\text{esc}}$, es decir,
cu\'anto desempe\~no pierde respecto de la referencia (un valor negativo
indica que mejora).

\begin{table}[!htb]
\centering
\caption[Escenarios de evaluaci\'on]{Condici\'on de referencia y $21$
escenarios de generalizaci\'on. Los escenarios con \texttt{extrap} en el
nombre superan las cotas del entorno ($S \le 9$, $T \le 10$).}
\label{tab:scenarios}
\small
\begin{tabular}{@{}llp{0.52\linewidth}@{}}
\toprule
\textbf{Familia} & \textbf{Identificador} & \textbf{Descripci\'on} \\
\midrule
referencia    & \texttt{sev\_base}             & Distribuci\'on de entrenamiento.\\
\midrule
severidad     & \texttt{sev\_uniform}          & $U\{0, \ldots, 9\}$, sin estructura modal.\\
severidad     & \texttt{sev\_gauss\_low}       & $\mathcal{N}(2, 1{,}5)$ truncada, brotes leves.\\
severidad     & \texttt{sev\_gauss\_mid}       & $\mathcal{N}(4{,}5, 2{,}0)$ truncada, brotes medios.\\
severidad     & \texttt{sev\_gauss\_high}      & $\mathcal{N}(7, 1{,}5)$ truncada, brotes intensos.\\
severidad     & \texttt{sev\_weibull}          & Weibull($k=1{,}5$), cola superior pesada.\\
severidad     & \texttt{sev\_beta\_lowskew}    & $9 \cdot \mathrm{Beta}(2, 5)$, sesgo bajo.\\
severidad     & \texttt{sev\_beta\_highskew}   & $9 \cdot \mathrm{Beta}(5, 2)$, sesgo alto.\\
severidad     & \texttt{sev\_bimodal}          & $0{,}5\,\mathcal{N}(2, 1) + 0{,}5\,\mathcal{N}(7, 1)$.\\
severidad     & \texttt{sev\_extrapolate\_high}& $U\{10, 11, 12\}$, fuera de rango.\\
\midrule
longitud      & \texttt{len\_all\_short}       & Todas las secuencias de longitud $3$.\\
longitud      & \texttt{len\_all\_long}        & Todas las secuencias de longitud $10$.\\
longitud      & \texttt{len\_geometric}        & $\mathrm{Geom}(p = 0{,}2)$, recortada a $[3, 10]$.\\
longitud      & \texttt{len\_poisson}          & $\mathrm{Poisson}(\lambda = 5)$, recortada a $[3, 10]$.\\
longitud      & \texttt{len\_extrapolate\_long}& $U\{11, \ldots, 20\}$, fuera de rango.\\
\midrule
conjunta      & \texttt{joint\_high\_long}     & Gaussiana alta $\times$ todas largas.\\
conjunta      & \texttt{joint\_low\_short}     & Gaussiana baja $\times$ todas cortas.\\
conjunta      & \texttt{joint\_uniform\_geom}  & Uniforme $\times$ geom\'etrica.\\
conjunta      & \texttt{joint\_extrap\_both}   & Severidad y longitud fuera de rango.\\
\midrule
estructural   & \texttt{struct\_few\_long\_blocks}  & $4$ bloques $\times$ $16$ secuencias.\\
estructural   & \texttt{struct\_many\_short\_blocks}& $16$ bloques $\times$ $4$ secuencias.\\
estructural   & \texttt{struct\_more\_total}        & $8$ bloques $\times$ $16$ secuencias ($n = 128$).\\
\bottomrule
\end{tabular}
\end{table}

\subsection{Protocolo de evaluaci\'on}

\begin{enumerate}
\item \textbf{Optimizaci\'on y reentrenamiento.} Cada modelo optimizable
se ajusta con Optuna y se reentrena desde cero con los mejores
hiperpar\'ametros. Q-Learning base se entrena con sus valores fijos.
\item \textbf{Evaluaci\'on determinista.} Cada modelo se eval\'ua en los
$22$ escenarios (Ap\'endice~\ref{ap:orchestrator}) eligiendo siempre la acci\'on factible de mayor valor
($\varepsilon = 0$). Los modelos individuales y los ensambles se
eval\'uan como dos grupos de seis modelos cada uno; Q-Learning base se
agrega como referencia.
\item \textbf{Comparaci\'on estad\'istica} (Secci\'on~\ref{sec:stats-bg}). Para cada modelo y escenario se calculan la media y la desviaci\'on est\'andar de $\bar r$, la degradaci\'on, y el $p$ de Welch y la $d$ de Cohen frente a su propia referencia. Tambi\'en se calcula la divergencia KL entre su distribuci\'on de acciones y la de la referencia. Para cada par de modelos de un mismo grupo se calculan
la $d$ de Cohen, el $p$ de Welch y la KL simetrizada entre sus
histogramas de desempe\~no, juntando las secuencias de los $21$
escenarios de generalizaci\'on. Como todos los
modelos enfrentan las mismas secuencias, la prueba de Welch se usa como
contraste de referencia; un contraste estrictamente pareado requerir\'ia
analizar las diferencias secuencia por secuencia.
\end{enumerate}

\biblio % Needed for referencing to working when compiling individual subfiles - Do not remove
\end{document}
